\documentclass[
  aps,
  reprint,
  superscriptaddress,
  floatfix,
  nofootinbib,
  longbibliography
]{revtex4-2}

\usepackage[T1]{fontenc}
\usepackage[utf8]{inputenc}
\usepackage{amsmath,amssymb}
\usepackage{graphicx}
\usepackage{bm}
\usepackage{xcolor}
\usepackage{hyperref}

\hypersetup{
  colorlinks=true,
  linkcolor=blue,
  citecolor=blue,
  urlcolor=blue
}

\begin{document}

\title{Imagers, \emph{Rechners}, and the Next Humiliation of Theory}

\author{Karl Svozil}
\email{karl.svozil@tuwien.ac.at}
\affiliation{Institute for Theoretical Physics, TU Wien,
Wiedner Hauptstra\ss{}e 8--10/136, 1040 Vienna, Austria}

\date{\today}

\begin{abstract}
Theory is not one-shot proof production but a traffic between image and
formalism. Current language models first enter this traffic as artificial
\emph{Rechners}: formalizers, programmers, critics, and prose machines. This
temporarily favors the visionary researcher, whose bottleneck is often closure,
while unsettling the human \emph{Rechner}, whose prestige rested on scarce formal
fluency. The division is unstable: with stochastic exploration, memory,
self-criticism, tools, and multimodal grounding, machines may move from closing
images to generating them. The relevant object is therefore the
human--model--tool--memory loop, not the isolated model.
\end{abstract}

\maketitle

\section{The wrong benchmark}

The \textit{First Proof} experiment, reported for a wider public by Roberts,
should be granted its narrow result~\cite{roberts2026,firstproof}. Present public
systems, asked to solve hard unpublished mathematical problems in a single
autonomous pass, are unreliable. They drift, hallucinate, solve nearby problems,
smuggle in conclusions, and write with the confidence of a mediocre professor
after dinner. That is worth knowing. It is not the central fact about theory.

A theorist does not receive a problem from heaven, produce a proof in one
exhalation, and then walk to lunch. A theorist broods, sketches, changes
notation, makes a wrong analogy, rescues a fragment, checks a sign, loses a
boundary term, consults a colleague, rewrites the statement, weakens the theorem,
and only then obtains something that may deserve print. Theory is not an answer.
It is a recurrent machine.

The question, then, is not whether a language model is already an autonomous
mathematician. It is whether a loop consisting of human taste, several models,
symbolic tools, numerical checks, and persistent memory enlarges the space of
ideas that can be made explicit, criticized, and repaired. In daily practice the
answer is already yes. The model is not a small Einstein in a box. At present it
is closer to an artificial \emph{Rechner}: a formalizing hand that can turn an
image into a lemma, a diagram into code, or a hunch into a claim that can finally
be broken. The phrase is useful only if one also sees its expiration date. A
machine that begins as a formalizer need not remain there.

\section{Images need hands}

Theoretical physics has always been a conspiracy between image and formalism.
Hertz made this explicit: theories are inner images whose necessary consequences
must correspond to the necessary consequences of the things represented~\cite{hertz-94,lutzen2005}.
Images are also at the core of research programs: they supply the organizing
Ansatz, the heuristic nucleus around which formal development, criticism, and
repair can proceed~\cite{lakatosch}. The image comes first: a gear train, a
light ray, a curved ruler, a forbidden coloring, a rotating universe, a bridge
over an abyss. But an image is still only a temptation. It becomes physics when
it survives algebra, units, boundary conditions, limiting cases, signs, and
hostile readers.

The distinction is deliberately schematic. It is not meant as a psychology of
genius, still less as a claim that image and formalism are separable substances.
In serious theory they often co-produce one another: the algebra may force a new
image into view, and a good image is already partly formal. I use the opposition
more modestly, as a temporary coordinate system for the present technological
moment. Current language models are stronger at producing formal variants, code,
notation, criticism, and prose than at originating stable scientific images; the
distinction helps name that asymmetry without pretending that it is
metaphysically deep.

The history of physics is full of this double act, although heroic memory prefers
solitary faces. Einstein's relation to the mathematical environment around
general relativity, the Hilbert--Einstein priority dispute, and the later unified
field work with Walther Mayer should not be flattened into a childish story of
genius and servant~\cite{corry1997,pais,lessel2023}. Yet the word
\emph{Rechner}, attached to Mayer in recollection, is too revealing to discard.
It names an indispensable competence while diminishing it: here is the hand that
closes the image; do not mistake the hand for the seer.

That hierarchy was never secure. Schr\"odinger's physical Ansatz for hydrogen did
not by itself close the calculation. The radial equation had to be recognized and
controlled as a special-function problem, with Laguerre polynomials and a
termination condition yielding the spectrum; biographical accounts connect part
of this closure with Weyl~\cite{ANDP:ANDP19263840404,Moore-Schroedinger}. The
Einstein--Podolsky--Rosen paper gives the opposite warning: formal apparatus can
also bury the point. Einstein complained to Schr\"odinger that the essential
matter had been smothered by erudition~\cite{epr,fine-sh}. A formalizer can save
an image. A formalizer can also suffocate it.

This history explains why language models matter first to one temperament rather
than another. If one thinks in images, analogies, possible mechanisms,
half-formed geometries, and dangerous sentences, then the bottleneck is often not
invention but closure. The algebraic hand, the code, the notation audit, and the
hostile reader are missing. A present language model supplies much of that hand,
cheaply and repeatedly. For the visionary researcher this is not replacement. It
is prosthesis.

For the human \emph{Rechner}, the same fact is less flattering. If one's special
power is fast derivation, clean notation, proof criticism, literature digestion,
and polished exposition, then the artificial hand enters precisely the old
prestige niche. It is not generally better; it is often foolish. But it is fast,
tireless, improving, and rentable as compute. The humiliation begins not when the
machine is always right, but when it is sometimes right enough.

\section{Good news for dreamers, bad news for calculators}

The present technology currently favors the ``artsy'' researcher: the one with
too many images and not enough formal patience. Such a researcher can now say:
here is the picture; make it a lemma; attack the lemma; write the code; find the
sign error; render the figure; compress the argument; now tell me why it is
false. This is not magic. It is the externalization of a scarce collaborator.

The migration came in two steps. Around 2025 it became natural to stop writing
much numerical and symbolic code by hand and instead specify the computation:
produce a Python check, a Mathematica reduction, a finite search, a TikZ diagram.
The author moved from typing syntax to specifying intent. From 2026 onward the
same move reached the manuscript. The article itself became specifiable: here is
the argument, here are the images, here is the desired wound, here are the
historical anchors, here is the level of provocation; now draft it, attack it,
shorten it, restore the sting. One may object that the words are no longer purely
the author's. Correct, but uninteresting. The harder question is whether the
author owns the images, the cuts, the corrections, the responsibility, and the
right to reject normalized mush.

Technique therefore shifts from an identity to an orchestration problem. Formal
competence does not vanish; it becomes necessary for judgment. But it loses some
scarcity value. Prestige moves upstream: toward problem choice, image formation,
taste, severity of tests, recognition of empty fluency, and responsibility for
what survives.

Again, the distinction must not be overdrawn. Not every kind of formalization
has been automated to the same degree. There is a shallow end: clean prose,
\LaTeX, code scaffolding, notation checks, and referee-style objection lists.
There is also a deep end: recognizing the right mathematical structure, choosing
the correct abstraction, seeing that a problem is secretly spectral, geometric,
variational, cohomological, or combinatorial. Present models are much better at
the first than at the second. The claim is therefore not that Mayer, Weyl, or a
first-rate mathematical physicist has been replaced by autocomplete. It is that
a growing part of the surrounding formal labor has become cheap enough to alter
the division of work and prestige.

This advantage is tactical, not metaphysical. The old consolation was that
machines execute while humans imagine. For now, models are more reliable as
closers than as originators of deep images. But stochastic exploration,
higher-temperature sampling, long-horizon memory, retrieval, proof tools,
self-criticism, and multimodal grounding are exactly the ingredients one would
add if one wanted machines to become less clerical and more visionary.

\section{The loop is the machine}

A single model call is almost never the relevant object. The relevant object is
the loop: one model proposes, another attacks, a symbolic or numerical tool
checks, the human adjudicates, the draft remembers, and the next pass starts from
the scar. This is not a metaphor but a computational architecture. Turing's
lesson is not that intelligence is a clever bounded answer; it is that
computation is an extended process with state~\cite{turing-36}. In research, the
state consists of failed routes, conventions, lemmas, sign errors, referee
objections, accepted definitions, and old embarrassments.

The working pattern is simple. Begin with an image or conjecture. Ask for a
formalization. Ask another pass to destroy it. Check the disputed steps by code,
dimensions, limiting cases, known examples, or proof assistants where available.
Keep unresolved objections instead of smoothing them away. Then assign status:
proved, plausible, speculative, false, or not yet worth naming. The model is not
an author in the responsible sense. It is a source of formal pressure.

This pressure can also pollute. If formal-looking arguments become cheap while
verification remains expensive, the result is not enlightenment but a fog of
plausible wrongness. The loop is valuable only when it includes adversarial
checking, external tools, and a human who is competent enough to reject fluent
nonsense. An artsy researcher without any formal discipline is not liberated by
the machine; he is merely supplied with better-looking errors.

Nor should one pretend that the human remains sovereign in some childish sense.
The prompt is not nothing, but it is not everything. A principal who asks a
question and judges an answer has not thereby performed every decisive
transformation inside the answer. Conversely, the formalizer without the image may
optimize the wrong thing magnificently. The honest description is not command and
obedience but unstable coupling.

\section{The next humiliation}

There is a professional injury here. Freud's sequence of humiliations---cosmological,
biological, psychological---has often been extended to the digital
case~\cite{freud1917,floridi2014}. The present wound is more local: it strikes
the prestige economy of formal theory.

The mathematically trained theorist has long possessed a secure currency:
notation, derivation, clean exposition, proof criticism, the ability to convert
fog into accountable mathematics. A frontier model can now produce a plausible
derivation, a referee-style critique, a code translation, a notation audit, and a
polished paragraph in seconds. Often it is wrong. Sometimes it is stupid. But the
insult is that wrongness is no longer the whole story.

This is not an \emph{ad hominem} argument against formal criticism. If a model
finds a missing hypothesis, the hypothesis is missing; if a human formalist finds
a sign error, the sign is wrong. But psychology matters to institutions. A
community that built rank around scarce formal fluency must now face a device
that mass-produces formal fluency. The visionary feels enlarged because the
missing hand arrived. The calculator feels displaced because the hand was the
identity.

The response is not to despise technique. Without technique the visionary becomes
a fog machine. This is not the abolition of hierarchy but its relocation.
Prestige moves from the mere possession of formal fluency toward command over the
whole loop: choosing fertile images, designing severe tests, knowing when
formalization has killed the thought, and assuming responsibility for the final
claim. The new aristocracy, if there must be one, should not be the people who
can type the cleanest derivation. It should be the people who know which
derivation deserves to exist.

\section{Will machines become visionary?}

The tempting line is that humans have images and machines have formalization. For
the present generation this is often pragmatically adequate. It is not a theorem.

Human images themselves may arise from processes less dignified than our
self-descriptions: recombination, analogy, dream, error, obsession, noise, bodily
state, music, landscape, and long unconscious sorting. Poincar\'e and Hadamard
emphasize the alternation between unconscious variation and conscious selection
in mathematical invention~\cite{poincare1908science,hadamard1945psychology}. That
alternation is not alien to machine design. A system that samples widely,
reflects on its failures, maintains memory over months, retrieves across domains,
uses tools, and tests its conjectures would no longer be merely a clerk. It would
begin to approximate, functionally if not phenomenologically, the ecology from
which human images often emerge.

Temperature is a crude symbol for this issue. Too low, and the model becomes a
civil servant of the average sentence. Too high, and it becomes delirium with
grammar. Creativity has always required an alliance between variation and
selection: randomness proposes, criticism disposes. If future systems learn to
manage that loop internally, the old division between visionary and
\emph{Rechner} will erode.

Multimodal grounding may matter as well. Many theorists do not report ideas as
sentences. They arrive as diagrams, motions, tensions, rhythms, impossible
objects, or spatial compressions. Music, walking, desert, mountain, and shoreline
are not mystical decorations; they are structured nonverbal regimes in which
symmetry, expectation, interruption, and resolution are felt before they are
propositional. A model acting only in text may imitate the report of such states.
A model moving among visual, symbolic, auditory, and kinesthetic structure may
generate stranger and less clerical images.

The artificial \emph{Rechner} may remain a \emph{Rechner}; it may also become a
competitor in imagination. The tactical advantage of the artsy researcher is
that, today, the machine supplies the missing hand. The strategic risk is that
tomorrow it may also grow eyes.

\section{The metaphysical escape hatch}

There is a deeper human defense---dualism: perhaps human minds are interfaces to
something beyond the material world. Perhaps they do not merely recombine,
sample, and select, but touch truth itself: Platonistic, religious, dualistic,
or ineffable in the sense analyzed by Jonas~\cite{Jonas-ineffability}. This
should not be mocked. G\"odel defended mathematical realism and a form of
intuition not reducible to formal syntax~\cite{godel1964cantor,parsons1995platonism},
and Ramanujan remains the disturbing example of formulaic vision outrunning
proof~\cite{hardy1940ramanujan,berndt1998ramanujan,berndt2020living}.

For science the rule is brutal. The flash must enter public space. A theorem
seen inwardly is not yet mathematics. A physical image is not yet physics. If it
can be formalized, checked, attacked, and repaired, the \emph{Rechner} enters
again. If it cannot, it may be profound, but it remains private. The machine
need not possess the ecstasy to participate in the domestication of its
consequences.

\section{The prober enters the archive}

The machinery will not remain politely inside private drafts. It will enter peer
review, literature search, replication, and accusation. A model that can find the
weak step in a new manuscript can also be pointed at a paper from 1978. This is
the software-fuzzer analogy: the stronger prober does not create the bug, but it
makes the bug cheap to find.

That prospect is hygienic and terrifying. A machine-assisted scan of the
scientific record could expose errors, unsupported inferences, duplicated
figures, citation fraud, and hidden gaps at a scale no human referee system could
match. It could also industrialize suspicion. Honest mistakes could become
accusation fodder; authors could be buried under unweighted objection lists;
referees could outsource paranoia while keeping the authority.

The sane norm is responsibility. A machine flag is not a verdict. A
machine-generated objection is not a referee report until a human adopts it,
weights it, and becomes answerable for it. Confidential manuscripts must not be
fed into unapproved systems. Authors should not have to answer every speculative
objection emitted by a model. But pretending that machine-assisted review will
not happen is childish. The prober has arrived.

\section{What should be tested}

The natural complement to one-shot benchmarks is not an easier benchmark but an
interactive one. Call it \emph{RechnerBench}, with irony: the test should detect
when the \emph{Rechner} becomes more than a \emph{Rechner}. Give the same hard
problems to several regimes: a one-shot model, an interactive single model, a
multi-model adversarial loop, a tool-grounded loop, a human alone, and a human
with the full loop. Measure not only final correctness, but auditability, false
starts, citation accuracy, time saved, compute cost, error detection, novelty,
and final-artifact quality.

Such a test would not flatter the machine. It would test the thing actually
being used. If the human--model--tool loop is the new research object, then
scoring the isolated model is like judging a laboratory by interviewing one
pipette.

The aim would be to reach, for theory, what Go became for strategic search:
a compact arena in which simple rules,
iteration, and feedback reveal capacities not visible in a single move. AlphaGo
mattered not only because it won, but because some of its moves first looked
alien and only later became intelligible as strategy~\cite{Silver-2017}.
A serious test of machine-assisted theory should likewise ask whether iterated
human--model--tool systems can produce surprising, auditable structures that no
isolated participant would have produced in one pass.



\section{Conclusion: the hand, the eye, and the wound}

The present language model is not an oracle and not a colleague in the full moral
sense. It is first an artificial \emph{Rechner}: a formalizing power inserted
into the old traffic between image and calculation. This helps the visionary
researcher because it supplies a missing hand. It wounds the human
\emph{Rechner} because it externalizes part of a formerly scarce identity. And it
may not stop there.

The human advantage, for now, lies in image, taste, responsibility, stubbornness,
and the ability to know that a smooth paragraph has betrayed the thought. The
machine advantage lies in relentless formal variation, cheap criticism, and
increasingly persistent memory. Perhaps humans have access to an ontological
depth machines lack. Perhaps future machines will learn to generate images as
well as close them. Neither possibility should be settled by slogans. The scandal
is already sufficient: the artificial \emph{Rechner} has arrived, the visionary
has been temporarily empowered, the calculator has been professionally unsettled,
and the old arrangement of credit, prestige, and humiliation no longer works.

\begin{acknowledgments}
The author acknowledges the use of frontier language models in programming,
drafting, compression, stylistic revision, and adversarial criticism of this
manuscript. The point is thematic: the paper describes a migration from
hand-written code to specified computation, and from hand-written prose to
specified, machine-drafted, human-adjudicated text. The models were not used as
autonomous factual authorities. The author assumes full responsibility for all
claims, errors, and interpretations. No AI system is listed as an author because
authorship requires responsibility and accountability. The author acknowledges
support by the TU Wien Library Open Access Funding Programme.
\end{acknowledgments}

\bibliography{svozil}

\end{document}
